\section{Conclusion}
\label{sec:conclusion}

We replicated the FSE~2026 posterior-guarantee study of configuration
sampling strategies across six corpora and 851 decision scopes, and
re-indexed its decisions by measured acquisition cost. The framework
replicates without qualification: posterior risk and confidence
discriminate strategies in $\approx$95\% of scopes, on targets from
encoding time to kernel size to model accuracy, with no learner in the
loop. The recommendations attached to it are population-relative: \srs's
superiority holds at the original study's budget scale on configurable
software and vanishes elsewhere; under measured cost it is never the best
choice outside the replication corpora, though often a defensible one at
a quantifiable premium.

The cost re-indexing yields the paper's central empirical fact---the
guarantee/cost Pareto frontier differs from the sample-size frontier in
97.5\% of scopes---and the equal-spend experiments explain the
recommendation instability mechanistically: cost--target coupling
determines whether cheap volume can substitute for expensive information,
and budget scale determines whether guided search amortizes. A counted
sample is not a faithful budget for guarantees; and ``cost-effective'' is
not a property a sampling strategy has, but one it exhibits relative to
an outcome and a coupling. We release all code, result tables, and a
command-level manifest to make both the replication and the extension
independently checkable.
